\section{Quadro comparativo}

{\footnotesize
\begin{longtable}{@{}L{2.1cm}L{4.2cm}L{4.2cm}L{4.2cm}@{}}
\toprule
& \textbf{Ahmed e Wahed (2020)} & \textbf{McElheran et al. (2024)} & \textbf{Knez (2023)} \\
\midrule
\endhead
\bottomrule
\endlastfoot
\textbf{Problema} &
A pesquisa de ponta em IA se concentrou em poucas organizações depois de 2012? Por falta de
acesso a computação? &
Quem usa IA nas empresas dos EUA, que tecnologias e onde? Faltava uma medida representativa. &
Por que tecnologias chegam tarde aos países pobres? Os modelos de adoção ignoram o
desenvolvimento desigual. \\ \addlinespace
\textbf{Dados} &
171.394 artigos de 57 conferências de computação, 2000--2019 (Scopus). Afiliações classificadas
por Fortune 500 e QS. &
ABS 2018 (ano-base 2017) ligado à LBD: 447 mil empresas (4 milhões ponderadas); 75 mil
\emph{startups}. &
CHAT (36 tecnologias, mais de 161 países, cerca de 200 anos) e Penn World Table (salários);
32.552 observações. \\ \addlinespace
\textbf{Método} &
Controle sintético generalizado, com o ImageNet 2012 como choque. Regressão de efeitos fixos com
série de \emph{compute}. TF-IDF dos resumos. &
Descritivas ponderadas e modelos de probabilidade linear (MQO) com controles de estado por
indústria. Corte transversal. &
Modelo de difusão (equação de Fick) no espaço de custo relativo. Calibração por MQ não lineares e
efeitos mistos. \\ \addlinespace
\textbf{Baseline} &
Contrafactual sintético a partir de 47 conferências de fora da IA. Placebos de área e de data. &
Empresas que não usam IA. Não há modelo só com controles. &
Nenhum modelo concorrente. A logística é citada, mas não estimada. \\ \addlinespace
\textbf{Métrica} &
ATT em artigos por conferência por ano, com IC por \emph{bootstrap}. &
Probabilidade de uso (pontos percentuais), taxas ponderadas por empresa e por emprego, $R^2$. &
Parâmetros $D$ e $\alpha$, $R^2$ (só no MQ não linear), curvas simuladas. \\ \addlinespace
\textbf{Resultado} &
Grandes empresas de tecnologia +44,5 e QS 1--50 +40,2; QS 201--300 $-$14,3 e QS 301--500 $-$5,6
(Tabelas~2 e~4). &
5,8\% das empresas usam IA (18,2\% por emprego). Associação com capital de risco, inovação de
processo e propriedade intelectual. Efeito do local some com controles. &
$R^2 = 0{,}474$. Países com 25\% do salário médio ficam décadas perto de zero na densidade
relativa. \\ \addlinespace
\textbf{Limitação} &
\emph{Compute} não é medido por organização. Talento, dados e \emph{compute} não são separados.
\emph{Preprint}. &
Uma única fotografia de 2017, antes da IA generativa. Uso autodeclarado com limiar baixo. Só os
EUA. &
Salário confunde custo com desenvolvimento. Figuras simuladas, sem dados. Sem código. \\
\end{longtable}
}
